%=====================================================================
\chapter{Storage Virtualization}\label{ch:storage}
%=====================================================================
\locator{4}{Part IV --- Storage, Network and Cloud Virtualization}

\begin{outcomes}
By the end of this chapter you will be able to:
\begin{tight}
  \item \textbf{Distinguish} block, file and object storage by the interface
        each presents, and \textbf{say} which a virtual machine needs.
  \item \textbf{Explain} what a virtual disk file is, and \textbf{compare} raw
        and qcow2 by what each gives up.
  \item \textbf{Compute} a thin-provisioning overcommit ratio and
        \textbf{predict} when a pool will exhaust.
  \item \textbf{Describe} copy-on-write snapshots and \textbf{explain} why a
        long snapshot chain degrades reads.
  \item \textbf{Explain} how software-defined storage places replicas, and
        \textbf{state} what it replaces.
\end{tight}
\end{outcomes}

\begin{prereq}
\chref{paravirt} (virtio-blk, and why an emulated disk controller costs 80\% of
its throughput) and \chref{what} (encapsulation --- a running machine as a
file, which this chapter makes literal). From Operating Systems you need
filesystems, blocks and the page cache.
\end{prereq}

\begin{keyterms}
block storage $\bullet$ file storage $\bullet$ object storage $\bullet$ LUN
$\bullet$ SAN $\bullet$ NAS $\bullet$ iSCSI $\bullet$ datastore $\bullet$
virtual disk $\bullet$ raw $\bullet$ qcow2 $\bullet$ VMDK $\bullet$ VHDX
$\bullet$ cluster $\bullet$ L1/L2 table $\bullet$ thick provisioning $\bullet$
thin provisioning $\bullet$ over-subscription $\bullet$ copy-on-write $\bullet$
snapshot chain $\bullet$ backing file $\bullet$ read amplification $\bullet$
software-defined storage $\bullet$ CRUSH $\bullet$ replica $\bullet$ erasure
coding $\bullet$ IOPS $\bullet$ discard / TRIM
\end{keyterms}

\section{Three Interfaces}

``Storage'' names three different contracts, and a great deal of confusion in
this area comes from using one word for all of them.

\rowcolors{2}{white}{lightbg}
\begin{longtable}{@{}L{2.4cm}L{4.4cm}L{3.9cm}L{4.0cm}@{}}
\toprule
\rowcolor{primary!15}
& \textbf{What it presents} & \textbf{Spoken over} & \textbf{Used for} \\
\midrule
\endhead
\term{Block} & An array of fixed-size blocks, with no notion of a file. Whoever
mounts it puts a filesystem on it & SAS, Fibre Channel, \term{iSCSI}, NVMe-oF &
\textbf{Virtual disks.} A guest expects a disk, so it must be block \\
\term{File} & A directory tree with files, locking and permissions, shared by
many clients at once & NFS, SMB & Shared datastores, home directories, anything
several machines must read together \\
\term{Object} & Whole immutable objects in a flat namespace, addressed by key
over HTTP & S3 API, Swift & Images, backups, logs, anything written once and
read many times \\
\bottomrule
\end{longtable}

\begin{conceptbox}[Why a Guest Must Have Block Storage]
A guest operating system contains a filesystem driver, and a filesystem driver
issues block reads and writes. It cannot be given an object store, because it
has no idea what an HTTP PUT is, and giving it a file share only works if
something underneath turns that into blocks --- which is exactly what a virtual
disk on NFS does.

\smallskip
So the chain is always: guest filesystem $\rightarrow$ block device
$\rightarrow$ virtual disk $\rightarrow$ whatever the host has. The interesting
engineering is in the last two steps, and the guest never sees any of it.
\end{conceptbox}

\dsfig{%
\begin{tikzpicture}[font=\scriptsize,
  bx/.style={rounded corners=2pt, line width=0.8pt, align=center,
             minimum width=5.0cm, minimum height=0.56cm, inner sep=3pt,
             font=\scriptsize},
  g/.style={bx, draw=secondary, fill=secondary!8, text=secondary!80!black},
  h/.style={bx, draw=primary, fill=primary!10, text=primary},
  s/.style={bx, draw=goldhl, fill=goldhl!10, text=goldhl!45!black},
  p/.style={bx, draw=neutral, fill=neutral!15, text=neutral!40!black},
  ar/.style={-{Stealth[length=1.8mm]}, primary!55, line width=0.8pt},
  note/.style={font=\scriptsize\itshape, text=neutral, anchor=west, align=left}
]
  \node[g] (a) at (0,4.5) {guest application: \texttt{write("hello")}};
  \node[g] (b) at (0,3.7) {guest filesystem (ext4)};
  \node[g] (c) at (0,2.9) {guest block device \texttt{/dev/vda}};
  \node[h] (d) at (0,2.1) {virtio-blk front end / back end};
  \node[h] (e) at (0,1.3) {\textbf{the virtual disk file} \ \texttt{vm1.qcow2}};
  \node[h] (f) at (0,0.5) {host filesystem or block device};
  \node[s] (g) at (0,-0.3) {LUN, pool or distributed cluster};
  \node[p] (h) at (0,-1.1) {physical media};
  \foreach \x/\y in {a/b,b/c,c/d,d/e,e/f,f/g,g/h}{\draw[ar] (\x)--(\y);}

  \node[note] at (2.8,4.5) {knows nothing about any of this};
  \node[note] at (2.8,2.9) {``a disk'' --- the illusion ends here};
  \node[note] at (2.8,2.1) {\chref{paravirt}: one exit per batch};
  \node[note] at (2.8,1.3) {\textbf{a file.} This is the whole chapter,\\
                            and the encapsulation of \chref{what}};
  \node[note] at (2.8,-0.3) {thin provisioning, snapshots and\\
                             replication all happen here};

  \draw[decorate, decoration={brace, amplitude=4pt}, accent, line width=0.8pt]
        (-2.7,4.85) -- (-2.7,2.6);
  \node[font=\scriptsize\itshape, text=accent!70!black, anchor=east,
        align=right] at (-2.9,3.7) {everything above\\the line is\\unmodified};
\end{tikzpicture}}%
{The path of one write. Each layer below the virtual disk file can be replaced
without the guest noticing --- which is what makes migration, snapshots and
replication possible at all.}{storage-stack}

\section{The Virtual Disk}

A virtual disk is a file on the host that the guest sees as a block device.
Two formats cover almost everything.

\begin{definitionbox}[Raw and qcow2]
A \term{raw} disk is a byte-for-byte image: block $n$ of the guest is byte
$n \times 512$ of the file. There is no metadata and no translation, so it is
as fast as the host filesystem and can be mounted directly with
\texttt{losetup}.

\smallskip
\term{qcow2} --- QEMU copy-on-write, version 2 --- adds an indirection table.
The file is divided into \term{clusters} (64 KB by default) and a two-level
L1/L2 table maps guest offsets to cluster locations. That indirection buys
three things raw cannot do: allocate clusters only when written (thin
provisioning), keep internal snapshots, and reference a \term{backing file} so
that several disks share one unwritten base.

\smallskip
VMware's \term{VMDK} and Microsoft's \term{VHDX} are the same idea with
different metadata, and \texttt{qemu-img convert} moves between all of them.
\end{definitionbox}

\begin{alertbox}[Raw is faster, and the gap is smaller than people assume]
qcow2 must consult its L2 table before reaching data, and must allocate and
update metadata on first write to each cluster. On a freshly created image with
small random writes the penalty is real --- tens of per cent.

\smallskip
Two things shrink it. The L2 tables are cached, so a steady-state workload
mostly hits the cache; and \texttt{qemu-img create -o preallocation=metadata}
builds the tables up front, removing the first-write cost. With a modern qcow2
and a warm cache the difference on typical server workloads is usually in the
low single digits.

\smallskip
Use raw when you need every last IOPS and will handle snapshots at the layer
below. Use qcow2 when you want the features, which is most of the time.
\end{alertbox}

\section{Thin Provisioning, and How It Fails}

\term{Thick provisioning} allocates all of a disk's space when it is created: a
500 GB disk takes 500 GB immediately. \term{Thin provisioning} allocates only
what has been written, so that same disk might occupy 40 GB.

Thin provisioning is the storage equivalent of the memory overcommit in
\chref{cpumem}, and it has the same shape: usually a large saving, occasionally
a cliff.

\dsfig{%
\begin{tikzpicture}[font=\scriptsize,
  bar/.style={draw=neutral!55, line width=0.7pt},
  note/.style={font=\scriptsize\itshape, align=left, anchor=west}
]
  % --- the three quantities ---
  \foreach \y/\lab in {2.6/{provisioned \ \ 20 TB}, 1.7/{physical \ \ 10 TB},
                       0.8/{written \ \ 7.2 TB}}{
    \node[font=\scriptsize\bfseries, text=neutral!40!black, anchor=east]
         at (-0.2,\y) {\lab};}

  \draw[bar, fill=secondary!15] (0,2.4) rectangle (12.0,2.8);
  \node[font=\tiny, text=secondary!80!black] at (6.0,2.6)
       {what the 40 guests believe they have};

  \draw[bar, fill=primary!18] (0,1.5) rectangle (6.0,1.9);
  \node[font=\tiny, text=primary] at (3.0,1.7) {what the pool actually holds};

  \draw[bar, fill=greenhl!30] (0,0.6) rectangle (4.32,1.0);
  \node[font=\tiny, text=greenhl!40!black] at (2.16,0.8) {what they have used};
  \draw[bar, fill=goldhl!25] (4.32,0.6) rectangle (6.0,1.0);
  \node[font=\tiny, text=goldhl!40!black] at (5.16,0.8) {2.8 TB};

  \draw[accent, line width=1.2pt, dashed] (6.0,0.4) -- (6.0,2.95);
  \node[font=\scriptsize\bfseries, text=accent, anchor=south] at (6.0,3.0)
       {the cliff};

  \draw[-{Stealth[length=2mm]}, greenhl!70, line width=1pt]
        (4.32,0.35) -- (5.3,0.35);
  \node[font=\tiny, text=greenhl!45!black, anchor=west] at (5.35,0.35)
       {+600 GB a month $\Rightarrow$ 4.6 months};

  \node[note, text=accent!70!black] at (0,-0.5)
       {\textbf{At the dashed line the pool is full.} Guests do not get a
        warning: a write fails, the guest filesystem sees an I/O error and\\
        remounts itself read-only. Databases stop. Several guests do this at
        once, because they share the pool.};
\end{tikzpicture}}%
{Thin provisioning, drawn to scale for the worked example. The gap between the
second and third bars is the saving; the gap between the first and second is
the exposure.}{thin-provisioning}

\begin{worked}[When Does This Pool Run Out?]
A datastore of 10 TB serves 40 guests, each with a 500 GB thin-provisioned
disk. Measured: mean 180 GB written per guest, growing 15 GB per guest per
month. Decide whether to act, and when.

\smallskip
\textbf{Step 1 --- the three quantities.}
\[ \text{provisioned} = 40 \times 500 = 20{,}000 \text{ GB} = 20 \text{ TB} \]
\[ \text{physical} = 10 \text{ TB} \qquad
   \text{written} = 40 \times 180 = 7{,}200 \text{ GB} = 7.2 \text{ TB} \]

\smallskip
\textbf{Step 2 --- the over-subscription ratio.}
\[ \sigma = \frac{20}{10} = \mathbf{2.0} \]
Every guest could legitimately fill its disk, and the pool can honour half of
them.

\smallskip
\textbf{Step 3 --- current headroom.}
\[ 10 - 7.2 = 2.8 \text{ TB free, or } 72\% \text{ of the pool consumed} \]

\smallskip
\textbf{Step 4 --- time to exhaustion at the measured rate.}
\[ 40 \times 15 = 600 \text{ GB per month} \qquad
   \frac{2{,}800}{600} = \mathbf{4.6 \text{ months}} \]

\smallskip
\textbf{Step 5 --- but that is the optimistic figure.} Three things consume the
pool faster than the trend line:
\begin{tight}
  \item \textbf{Snapshots.} Each one holds the blocks its parent has since
        overwritten. A forgotten snapshot on a busy database can grow by its
        whole working set.
  \item \textbf{Deleted files do not return space.} A guest deleting 50 GB
        frees nothing in the pool unless \texttt{discard}/TRIM is enabled all
        the way down. Most estates discover this configuration only here.
  \item \textbf{Correlation.} A patch cycle writes to every guest in the same
        hour, so consumption arrives in steps, not a slope.
\end{tight}

\smallskip
\textbf{Step 6 --- what happens at 100\%.} Not a slowdown. The pool cannot
satisfy a write; the guest receives an I/O error; its filesystem remounts
read-only to protect itself. \emph{Many guests hit this within seconds of each
other}, because they share the pool --- so the failure is simultaneous and
estate-wide, and recovery requires space that by definition is not there.

\smallskip
\textbf{Step 7 --- the recommendation.} Act now, not in four months.
\begin{tight}
  \item Alert at 80\% of the pool (2 TB free), which at the measured rate is
        about six weeks away. Alerting at 95\% leaves no time to buy anything.
  \item Enable \texttt{discard} end to end and reclaim; on a four-year-old
        estate this commonly returns 10--20\%.
  \item Report any snapshot older than seven days, every day.
  \item Add capacity, or reduce $\sigma$ by shrinking the 500 GB disks nobody
        asked for.
\end{tight}
\end{worked}

\section{Snapshots and Copy-on-Write}

A snapshot freezes the disk as it is now. The implementation is the obvious
one: stop writing to the current file and make it read-only, then direct all
new writes to a fresh \term{overlay} file that records only what has changed.
A read checks the overlay first and falls back to the frozen file beneath.

This is \term{copy-on-write}, and it is the same mechanism as the overlayfs of
\chref{containers}, one layer down.

\begin{conceptbox}[Why a Long Chain Is Slow, and Why It Is Dangerous]
Each snapshot adds a layer. A read must check each overlay in turn until it
finds the block, so a chain of twelve snapshots can require twelve lookups
where one would do --- \term{read amplification}. Writes are worse: the first
write to a block must copy it up from wherever it currently lives.

\smallskip
The danger is worse than the latency. Every file in the chain is required: lose
or corrupt one and every snapshot above it is unreadable. A guest with a
six-month chain has six months of fragile dependency, and the usual
discovery is that nobody knows which snapshots are still needed.

\smallskip
\textbf{A snapshot is not a backup.} It lives on the same storage as the
original, so it does not survive the array failing, the pool filling, or the
ransomware that encrypted both. Keep chains short, delete snapshots on a
schedule, and keep real backups elsewhere. \chref{migration} covers what a real
backup is.
\end{conceptbox}

\section{Software-Defined Storage}

The traditional answer to shared storage was a \term{SAN}: a dedicated array,
dual controllers, a Fibre Channel fabric, and a price to match. The
software-defined answer is to put ordinary disks in the servers you already
have and let software make them one pool.

\term{Ceph} is the open example. Data is split into objects; each object's
placement is computed by \term{CRUSH}, a deterministic function of the object
name and a map of the cluster's failure domains. There is no lookup table and
no metadata server in the data path --- any client can compute where a block
lives, which is what lets it scale to thousands of nodes. VMware's \term{vSAN}
does the same job inside a vSphere cluster.

\begin{conceptbox}[Replication or Erasure Coding: a Real Choice]
Both survive failures; they trade capacity against CPU and latency.

\smallskip
\textbf{Three-way replication} keeps three copies. Usable capacity is a third
of raw; a read is served by one disk; recovery is a copy. Simple and fast.

\smallskip
\textbf{Erasure coding} --- say 4 data plus 2 parity --- stores 6 chunks for 4
chunks of data. Usable capacity is $4/6 = 67\%$ of raw against replication's
33\%, so a petabyte of raw yields 670 TB instead of 330 TB. The cost is that
every write computes parity and every degraded read reconstructs it, so
latency and CPU both rise.

\smallskip
The usual split: replication for virtual machine disks, where latency matters
and the data is active; erasure coding for backups, archives and object
storage, where capacity dominates and nothing is latency-sensitive. Putting a
busy database on erasure-coded storage is a common and expensive mistake.
\end{conceptbox}

\begin{pitfall}
\begin{tight}
  \item \textbf{Monitoring provisioned capacity instead of written capacity.}
        The first tells you nothing about when you run out. Alert on the pool's
        free space, at 80\%.
  \item \textbf{Treating a snapshot as a backup.} Same storage, same failure,
        same ransomware.
  \item \textbf{Leaving snapshots behind after an upgrade.} The chain grows, the
        pool fills and reads slow down. Report anything older than a week.
  \item \textbf{Assuming deleting files inside a guest frees pool space.} It
        does not without \texttt{discard} enabled end to end, which is off by
        default in many stacks.
  \item \textbf{Thin-provisioning the pool \emph{and} overcommitting memory
        \emph{and} oversubscribing CPU on the same host.} Each is defensible
        alone; together they fail at the same moment, because the event that
        triggers one triggers the others.
  \item \textbf{Erasure-coding a database.} Capacity efficiency bought with
        write latency, on the workload least able to pay it.
  \item \textbf{Running a guest's disk on a filesystem that is itself
        copy-on-write, with qcow2 on top.} Two layers of copy-on-write
        fragment each other. Pick one.
\end{tight}
\end{pitfall}

\begin{lab}[Thin Pools, Snapshots and the Cliff]
Builds a small thin pool with LVM, fills it, and watches a guest filesystem go
read-only --- the worked example's step 6, at a safe scale. Expect forty
minutes.

\begin{lstlisting}[language=bash]
#!/usr/bin/env bash
# ch12_storage.sh -- thin provisioning, snapshots and exhaustion.
set -euo pipefail

# --- 1. Virtual disk formats: what they cost on disk. ------------------
cd /tmp
qemu-img create -f raw   plain.raw  10G
qemu-img create -f qcow2 thin.qcow2 10G
ls -lh plain.raw thin.qcow2          # apparent size
du -h  plain.raw thin.qcow2          # actual size: this is the point
qemu-img info thin.qcow2 | head -6

# --- 2. A backing file: many disks, one shared base. ------------------
qemu-img create -f qcow2 -b thin.qcow2 -F qcow2 child1.qcow2
qemu-img create -f qcow2 -b thin.qcow2 -F qcow2 child2.qcow2
du -h child1.qcow2 child2.qcow2      # nearly nothing until written
qemu-img info child1.qcow2 | grep -i backing

# --- 3. Build a thin pool, deliberately over-subscribed. --------------
sudo truncate -s 2G /tmp/pv.img
LOOP=$(sudo losetup --find --show /tmp/pv.img)
sudo pvcreate "$LOOP" ; sudo vgcreate vsvg "$LOOP"
sudo lvcreate -L 1.5G --thinpool vspool vsvg
# Four 1 GB volumes out of a 1.5 GB pool: sigma = 2.7
for i in 1 2 3 4; do
  sudo lvcreate -V 1G --thin -n "vol$i" vsvg/vspool
done
sudo lvs vsvg -o lv_name,lv_size,data_percent

# --- 4. Put filesystems on them and write a little. -------------------
for i in 1 2 3 4; do
  sudo mkfs.ext4 -q "/dev/vsvg/vol$i"
  sudo mkdir -p "/mnt/v$i" && sudo mount "/dev/vsvg/vol$i" "/mnt/v$i"
  sudo dd if=/dev/urandom of="/mnt/v$i/data" bs=1M count=200 status=none
done
sudo lvs vsvg -o lv_name,data_percent    # pool filling up

# --- 5. Now walk off the cliff. ---------------------------------------
# Keep writing until the pool cannot satisfy a write.
sudo dd if=/dev/urandom of=/mnt/v1/more bs=1M count=900 || \
  echo "write failed -- the pool is full"
dmesg | tail -5 | grep -iE 'thin|read-only|I/O error' || true
mount | grep -E '/mnt/v[0-9]'            # look for "ro" on several at once

# --- 6. Snapshots and read amplification. -----------------------------
cd /tmp
qemu-img create -f qcow2 base.qcow2 1G
sudo modprobe nbd max_part=8 || true
for n in $(seq 1 8); do
  qemu-img create -f qcow2 -b "$( [ $n -eq 1 ] && echo base.qcow2 \
    || echo "snap$((n-1)).qcow2" )" -F qcow2 "snap$n.qcow2"
done
qemu-img info --backing-chain snap8.qcow2 | grep -c '^image:'
# Eight files, every one of which is required to read the newest.

# --- 7. Reclaim: discard is what returns space. -----------------------
sudo fstrim -v /mnt/v2 || echo "discard not supported on this stack"
sudo lvs vsvg -o lv_name,data_percent

# --- 8. Tear down. ----------------------------------------------------
for i in 1 2 3 4; do sudo umount "/mnt/v$i" || true; done
sudo lvremove -f vsvg || true ; sudo vgremove -f vsvg || true
sudo pvremove -f "$LOOP" || true ; sudo losetup -d "$LOOP" || true
rm -f /tmp/{plain.raw,thin.qcow2,child*.qcow2,base.qcow2,snap*.qcow2,pv.img}
\end{lstlisting}

\textbf{What to notice.} Step 5 is the part worth sitting with. The failure is
not gradual and it is not confined to the volume that caused it: several
filesystems go read-only together, because the resource they were quietly
sharing ran out. That is the whole argument for alerting at 80\%.
\end{lab}

\begin{checkpoint}
\begin{tightnum}
  \item A pool of 40 TB serves 120 guests with 800 GB disks each, of which a
        mean 240 GB is written, growing 20 GB per guest per month. Give the
        over-subscription ratio, the headroom and the time to exhaustion. At
        what free-space figure would you alert, and why?
  \item A guest deletes 300 GB of files and the storage team reports no change
        in pool usage. Explain why, and give the setting that changes it.
  \item Why is a chain of twelve snapshots both slow to read and dangerous to
        keep?
\end{tightnum}
\end{checkpoint}

\begin{chaptersummary}
\begin{tight}
  \item Block, file and object are three different contracts. A guest needs
        block, because its filesystem driver issues block requests.
  \item A virtual disk is a file the guest sees as a device. raw is a
        byte-for-byte image; qcow2 adds an L1/L2 indirection that buys thin
        provisioning, snapshots and backing files.
  \item Thin provisioning is memory overcommit for disks: usually a large
        saving, occasionally a cliff. Monitor written capacity, never
        provisioned, and alert at 80\%.
  \item Pool exhaustion is not a slowdown. Writes fail, guest filesystems go
        read-only, and they do it together because they share the pool.
  \item Snapshots are copy-on-write overlays. Long chains amplify reads and
        make every file in the chain load-bearing. A snapshot is not a backup.
  \item Software-defined storage computes placement rather than looking it up.
        Replication for virtual disks, erasure coding for archives.
\end{tight}
\end{chaptersummary}

\begin{reviewq}
  \item \textbf{(MCQ)} A guest operating system requires its storage to be
        presented as: (a)~object storage; (b)~block storage; (c)~an NFS share;
        (d)~any of these.
  \item \textbf{(MCQ)} When a thin pool reaches 100\% of physical capacity:
        (a)~writes slow down; (b)~the oldest snapshots are deleted; (c)~writes
        fail and guest filesystems typically remount read-only; (d)~the
        hypervisor migrates guests away.
  \item \textbf{(Short)} State what qcow2's indirection table buys over raw,
        and what it costs.
  \item \textbf{(Short)} Explain read amplification in a snapshot chain, and
        why every file in the chain matters.
  \item \textbf{(Short)} Give the rule for choosing between replication and
        erasure coding, with the reason.
  \item \textbf{(Applied)} An estate thin-provisions storage at
        $\sigma = 2.4$, overcommits memory at 1.3 and oversubscribes CPU at
        2.0. Explain why these three are more dangerous together than
        separately, give a scenario that triggers all three at once, and say
        what you would change first.
  \item \textbf{(Applied)} You inherit an estate with 400 guests, 1{,}900
        snapshots and no snapshot policy. Write the plan you would follow to
        get to a safe state, in order, saying what you would measure before
        deleting anything.
\end{reviewq}
